% ===================================================================
%  TABELUL Nr. 15 — Piața. — Merindele.
%  Traduction 2026 d'apres texto/io, controlee sur texto/fr, texto/en
%  et texto/es. Consigne : voir texto/ro/10-tabelul-01.tex.
%
%  LE TABLEAU DU MARCHE EST LE PLUS TURC DE LA COLONNE APRES CELUI DE
%  LA MAISON, et c'est le meme critere : « telemea » et « cașcaval »
%  les fromages, « pastramă » la viande sechee, « ciorbă » la soupe,
%  « iaurt », « halva », « rahat », « ardei » le poivron, « vânătă »
%  l'aubergine, « dovleac » la courge. La cuisine s'apprend a table,
%  non a l'ecole, et elle prend ses mots ou elle prend ses plats.
% ===================================================================

%%K t15-apar-1 apar
\VUsaut{4.0mm}
\VUtitre{15.0pt}{60}{TABELUL Nr. 15}
\VUsaut{3.0mm}

%%K t15-tit-1 sub
\VUpk{11.4pt}{40}{Piața. --- Merindele.}
\VUsaut{3.0mm}

%%K t15-01-1 p
1. --- Cei mai mulți dintre negustorii care ne dau merindele trebuincioase sunt
strânși sub \VUgras{șopronul} \textsuperscript{(1)} \VUgras{halei} sau pe străzile
învecinate.

%%K t15-02-1 p
2. --- \VUgras{Mezelarul} \textsuperscript{(2)}, care vinde \VUgras{șuncă}
\textsuperscript{(3)}, \VUgras{sângerete} \textsuperscript{(4)},
\VUgras{cârnați} \textsuperscript{(5)}, \VUgras{crenvurști}
\textsuperscript{(6)} și \VUgras{slănină} \textsuperscript{(7)}, stă lângă
\VUgras{negustoreasa de unt și brânză} \textsuperscript{(8)}, a cărei tarabă este
plină de \VUgras{brânzeturi olandeze} \textsuperscript{(9)} cu coajă roșie, de
\VUgras{camemberturi} \textsuperscript{(10)}, de \VUgras{roți de gruyère} mari
\textsuperscript{(11)} și de \VUgras{bucăți de unt} proaspăt \textsuperscript{(12)}.

%%K t15-03-1 p
3. --- Înaintea ei \VUgras{zarzavagioaica} \textsuperscript{(13)} le oferă
mușteriilor \VUgras{roșii} frumoase \textsuperscript{(14)}, \VUgras{conopidă}
\textsuperscript{(15)}, \VUgras{salată} \textsuperscript{(16)}, \VUgras{varză}
\textsuperscript{(17)}, \VUgras{dovleci} \textsuperscript{(19)},
\VUgras{pepeni} \textsuperscript{(18)} și chiar \VUgras{ciuperci}
\textsuperscript{(20)}. Dar \VUgras{berarul}, care și-a oprit \VUgras{căruța de
livrare} \textsuperscript{(21)}, încărcată cu \VUgras{butoiașe de bere}
\textsuperscript{(22)}, chiar înaintea tarabei ei, îi împiedică pe mușterii să se
apropie. Grădinarul îi aduce un coș de \VUgras{anghinare} \textsuperscript{(23)}.

%%K t15-tit-2 sub
\VUsaut{4.0mm}
\VUpk{10.2pt}{40}{A vinde și a cumpăra.}
\VUsaut{3.0mm}

%%K t15-04-1 p
4. --- În piață este mare însuflețire, fiindcă a venit ceasul \VUgras{tocmelii}
\textsuperscript{(24)}. \VUgras{Gospodinele} \textsuperscript{(26)}, care vin aici
la cumpărături, se opresc pe drum înaintea tarabei \VUgras{negustoresei de
măruntaie} \textsuperscript{(25)}, ca să cumpere \VUgras{burtă} sau \VUgras{cap de
vițel}.

%%K t15-05-1 p
5. --- Un \VUgras{bucătar} gras \textsuperscript{(27)} se tocmește pentru un
\VUgras{ton} frumos cu \VUgras{peștăreasa} \textsuperscript{(28)}, care ridică
prețul cum îi vine. I s-au adus mulți \VUgras{țipari, calcani, păstrăvi} și
\VUgras{crapi}. \VUgras{Codul} ei și \VUgras{midiile} ei sunt foarte proaspete.

%%K t15-tit-3 sub
\VUsaut{4.0mm}
\VUpk{10.2pt}{40}{Ieftin și scump.}
\VUsaut{3.0mm}

%%K t15-06-1 p
6. --- \VUgras{Negustoreasa de păsări} \textsuperscript{(29)}, așezată lângă ea,
este foarte cinstită. Vânzând un \VUgras{curcan}, o \VUgras{gâscă}, o
\VUgras{rață} sau o simplă \VUgras{găină}, nu cere mai mult decât prețul drept.

%%K t15-07-1 p
7. --- În colțul pieței, la etajul întâi, a deschis de curând prăvălie o
\VUgras{lenjereasă} \textsuperscript{(30)}, la care mușteriii vor găsi totdeauna o
alegere foarte bună de \VUgras{fețe de masă}, de \VUgras{șervete, șorțuri} și de
\VUgras{ștergare de bucătărie}. La același etaj și-a așezat atelierul
\VUgras{cuțitarul} \textsuperscript{(31)}. Ascute \VUgras{cuțitele} negustoreselor
și le face grădinarilor \VUgras{cosoare} din \VUgras{oțelul} cel mai tare.

%%K t15-tit-4 sub
\VUsaut{4.0mm}
\VUpk{10.2pt}{40}{Băcănia.}
\VUsaut{3.0mm}

%%K t15-08-1 p
8. --- Sub atelierul lui s-a deschis nu demult o prăvălie mare de merinde. Nu mai
este simpla și modesta \VUgras{băcănie} \textsuperscript{(32)} de altădată, unde se
vindeau numai lucrurile de toate zilele --- \VUgras{ceai} \textsuperscript{(33)},
\VUgras{ciocolată} \textsuperscript{(34)}, \VUgras{căpățână de zahăr}
\textsuperscript{(37)} și \VUgras{săpun} \textsuperscript{(38)}. Aici se vinde
totul: \VUgras{biscuiți crocanți}, \VUgras{conserve} \textsuperscript{(35)},
\VUgras{lichioruri} \textsuperscript{(36)}, \VUgras{rom, coniac, kirsch, anason} și
\VUgras{curaçao}. Vânatul se găsește aici tot atât de lesne --- \VUgras{iepure}
\textsuperscript{(39)}, \VUgras{sturzi} \textsuperscript{(40)},
\VUgras{potârnichi} \textsuperscript{(41)}, \VUgras{prepelițe}
\textsuperscript{(42)}, \VUgras{rață sălbatică} \textsuperscript{(43)} sau
\VUgras{fazan} \textsuperscript{(44)} --- ca și roadele de la tropice, precum
\VUgras{bananele} \textsuperscript{(46)} și \VUgras{ananașii}.

%%K t15-09-1 p
9. --- \VUgras{Vânzătorul} foarte îndatoritor \textsuperscript{(45)} este gata să
dea tot ce i se cere: \VUgras{dulceață} \textsuperscript{(47)}, de coacăze sau de
gutui, \VUgras{untură} \textsuperscript{(48)}, \VUgras{castraveciori}
\textsuperscript{(49)} sau \VUgras{sardele} sărate \textsuperscript{(50)}.

%%K t15-09-2 p
Aceasta este foarte la îndemână pentru \VUgras{cumpărătoare}
\textsuperscript{(51)}, care găsesc lângă lucrurile cele mai obișnuite cele mai rare
timpurii și cele mai gustoase roade: \VUgras{pere} zemoase \textsuperscript{(52)},
\VUgras{prune} \textsuperscript{(53)}, \VUgras{struguri} \textsuperscript{(54)},
\VUgras{piersici} \textsuperscript{(55)}, \VUgras{migdale} \textsuperscript{(56)} și
\VUgras{nuci} \textsuperscript{(57)}.

%%K t15-tit-5 sub
\VUsaut{4.0mm}
\VUpk{10.2pt}{40}{Cumpărătorii.}
\VUsaut{3.0mm}

%%K t15-10-1 p
10. --- \VUgras{Orezul} \textsuperscript{(58)} și \VUgras{lintea}
\textsuperscript{(59)} stau lângă lada de \VUgras{scrumbii} afumate
\textsuperscript{(60)}; \VUgras{cartofii} \textsuperscript{(61)} și
\VUgras{fasolea} \textsuperscript{(63)} sunt vecine cu \VUgras{merele}
\textsuperscript{(62)}, cu \VUgras{smochinele} \textsuperscript{(64)}, cu
\VUgras{prunele uscate} \textsuperscript{(65)} și cu \VUgras{alunele}
\textsuperscript{(66)}. \VUgras{Ouă} proaspete \textsuperscript{(67)} și
\VUgras{măsline} \textsuperscript{(68)} sunt de asemenea în această prăvălie, așa
că \VUgras{coșurile} \textsuperscript{(70)} și \VUgras{plasele}
\textsuperscript{(73)} se umplu foarte repede.

%%K t15-11-1 p
11. --- \VUgras{Cafeaua} \textsuperscript{(72)} acestei case alese este pe drept
vestită printre \VUgras{slujnice} \textsuperscript{(69)} și bucătărese, fiindcă
bătrânul \VUgras{băcan} \textsuperscript{(71)} o prăjește el însuși cu cea mai mare
grijă.

%%K t15-tit-6 sub
\VUsaut{4.0mm}
\VUpk{10.2pt}{40}{Ce se poate vedea aici.}
\VUsaut{3.0mm}

%%K t15-12-1 p
12. --- Nu toți vin în piață după merinde. În mișcarea pitorească a uliței care duce
spre șopron vezi cel mai felurit norod. Lângă un \VUgras{lucrător de la abator}
bonom \textsuperscript{(75)}, care împinge înaintea sa o \VUgras{roabă}
\textsuperscript{(74)}, iată doi soldați în permisie, foarte mândri de uniformele
lor strălucitoare: un \VUgras{husar} \textsuperscript{(76)} în \VUgras{dolman}
albastru și cu \VUgras{căciulă cu penaj}, și un \VUgras{caporal de zuavi} cu
șalvari largi și cu \VUgras{fes} pe cap.

%%K t15-13-1 p
13. --- \VUgras{Brutarul} \textsuperscript{(80)}, care stă în pragul
\VUgras{brutăriei} \textsuperscript{(78)}, tocmai a frământat aluatul pentru
\VUgras{pâinea} sa \textsuperscript{(79)} și a scos din cuptor \VUgras{cornuri}
\textsuperscript{(81)}, \VUgras{chifle} \textsuperscript{(82)} și
\VUgras{pâini rotunde} \textsuperscript{(83)}, care stau în vitrină.

%%K t15-14-1 p
14. --- Deasupra brutăriei locuiește o \VUgras{croitoreasă} iscusită
\textsuperscript{(84)}, la care lucrează câteva \VUgras{brodeze}
\textsuperscript{(85)}. Lângă ea locuiește o \VUgras{spălătoreasă și călcătoreasă}
\textsuperscript{(86)}, care scrobește \VUgras{rufele} \textsuperscript{(88)}, după
ce le spală și le usucă, și le calcă cu \VUgras{fierul de călcat}
\textsuperscript{(87)}.

%%K t15-tit-7 sub
\VUsaut{4.0mm}
\VUpk{10.2pt}{40}{Carne, pește și zarzavat.}
\VUsaut{3.0mm}

%%K t15-15-1 p
15. --- În \VUgras{măcelăria} \textsuperscript{(89)} de la parter se vinde orice fel
de carne --- \VUgras{carne de miel} \textsuperscript{(90)}, \VUgras{carne de vită}
\textsuperscript{(94)} sau \VUgras{carne de vițel} \textsuperscript{(92)} --- sub
chip de \VUgras{pulpe de miel, mușchi, fripturi} și \VUgras{cotlete}.
\VUgras{Ajutorul măcelarului} \textsuperscript{(93)} taie toată această carne și
pune deoparte \VUgras{oasele cu măduvă} \textsuperscript{(96)}. La ușa măcelăriei
stă \VUgras{negustoreasa de stridii} \textsuperscript{(97)}, care vinde
\VUgras{stridii} \textsuperscript{(98)}. Le deschide, dacă i se cere.

%%K t15-16-1 p
16. --- Toți acești negustori plătesc patenta sau taxa de loc; iată de ce
\VUgras{sergentul de stradă} \textsuperscript{(100)} le fugărește fără milă pe
bietele \VUgras{precupețe de zarzavat} \textsuperscript{(99)}, care le fac
concurență, plimbând pe cărucioarele lor \VUgras{morcovi, ridichi, gulii, praz} sau
\VUgras{legături de sparanghel, mazăre verde, spanac, măcriș, creson} și
\VUgras{pătrunjel}.

%%K t15-tit-8 sub
\VUsaut{4.0mm}
\VUpk{10.2pt}{40}{Ferește-te de sergent!}
\VUsaut{3.0mm}

%%K t15-17-1 p
17. --- Băiatul care vinde \VUgras{usturoi} \textsuperscript{(101)} și
\VUgras{ceapă} știe prea bine că nici lui nu îi este îngăduit să își vândă marfa
lângă piață. O ia la fugă, primejduindu-se să răstoarne un coș de \VUgras{crabi},
de \VUgras{raci} și de \VUgras{creveți}, care este al \VUgras{negustoresei}
\textsuperscript{(102)} de \VUgras{languste} și de \VUgras{homari}.

%%K t15-18-1 p
18. --- Toți se agită și fac mare zgomot; dar aceasta nu împiedică să se audă
limpede strigătul \VUgras{geamgiului} ambulant \textsuperscript{(103)} și al
\VUgras{cărbunarului} \textsuperscript{(104)}, care tocmai și-a lăsat o clipă
căruciorul, ca să ducă un \VUgras{sac de cărbuni} \textsuperscript{(105)}.
\VUsaut{6.0mm}
\VUfilet{22mm}
